\clearpage
\section*{September 20, 2026: Development-ground rules, parent mergers, and adopted site areas}
\textit{Agent-drafted research entry.}

From September 16 to 22 we read recorded deeds, zoning declarations, ground
leases, agency site plans and city tax maps for parents whose land the
automatic parcel match could not confirm. The review asked two questions.
Which physical ground did each development occupy, and which separately filed
buildings form one economic parent? It began with the post-policy parents that
contain an exact-99 filing and then turned to large historical parents.
GPT-5.6 Sol agents gathered documents in bounded groups; the supervisor read
the decisive original pages before accepting any finding. Two early agent
proposals to set land from later map polygons were rejected, and a Broadway
zoning drawing known only from a search-index extract was not treated as
evidence. The supervisor-exported ACRIS PDFs have no text layer. A two-page
local OCR test recovered the needed passages, but the proposed Jev reading
pilot was not run, and no machine classification entered any decision.

This entry records the rules used, the links accepted or left open, and the
adopted ground for every reviewed parent. It replaces the September 16 and 20
site-review entries and the boundary parts of the September 20 and 22 batch
entries. Unit-source findings from those batches are recorded separately.

\paragraph{What counts as development ground.}
The framework holds the observed economic-parent boundary fixed across
regimes and measures historical characteristics inside it. Development ground
is the physical land under the parent's buildings plus documented common land
such as courtyards and accessory parking. Earlier floor is the building floor
area recorded on that ground in the named pre-filing MapPLUTO release; it
enters existing built FAR. Five rules governed the adjudication.

First, development rights, zoning lots and retained buildings are not
development ground. At Ocean Parkway the church conveyed 64,000 square feet of
residential development rights. Those rights permit more floor space but add
no land. The adopted ground is the recorded 200-by-100-foot rectangle, 20,000
square feet, carved from the church's 57,000-square-foot old lot 25; the
church's building and its retained 40,000 square feet are excluded. Production
had carried 3,000 square feet because the same tax-lot identifier described a
different parcel in the reference map. Zoning lots are often much larger than
the ground: GO Broome's zoning lot is 51,884 square feet against 32,395 of
ground, and Beach 30th's City Planning plan describes a
63,346.60-square-foot zoning lot for two buildings that includes parcels with
other live filings. Easements and
access rights add no ground either, including Broadway's adjacent lot 28 and
East 108th Street's retained lot 342, which provides parking and access. Floor
on a retained neighboring building is removed from earlier floor; at Bedford
Square this is the 175,875-square-foot Sears building.

Second, a ground lease contributes the ground it assigns, not the lessor's
whole parcel. Casa Celina's executed June 2021 lease (ACRIS 2021063001174002,
PDF p.~9) bounds approximately 12,613.61 square feet: a
13,046.67-square-foot chord polygon less a 433.06-square-foot circular
segment. The adopted successor area of 12,613 replaces a 344,900-square-foot
campus value. A later NYCHA budget document says 9,400 square feet without a
corresponding boundary; that discrepancy is recorded, and the executed lease
governs. Penn South's lease bounds a 247-by-163-foot parcel inside a
393,100-square-foot campus. North 8th Street's 99-unit ground lease covers lot
133 and leaves the church on the neighboring parcels. The Flatbush tower's
ground lease replaces a 1,198-square-foot source lot.

Third, later documents can reconstruct a boundary, but they do not prove
ownership or the boundary at filing. The framework does not require proof that
the same subdivision or ownership existed on the reference date. The
\texttt{area\_basis} field records which measurement is used. Starhill's
combined ground is a later administrative measurement of the boundary drawn in
its March 2022 declaration, not a February 2020 survey. The 1580 Story Avenue
declaration establishes observed ground, not title at the June 2019 filing.
Zoning and FAR always come from the named pre-filing release, and several
earlier parcels share one decision row only when they share both FAR measures.

Fourth, a later vacant parcel does not establish earlier vacancy. Casa
Celina's later vacant-land code follows site preparation, so its archived
prior-use category stays; De Kalb keeps its old public-facility category.
Dated evidence can still correct the archive. City Planning dates complete
demolition at GO Broome to June 2019, before the November 2019 reference
snapshot. Following Jacob's instruction on September 22, the stale
4,600-square-foot archived floor is deducted as a documented source
correction, leaving zero earlier floor. Its ground and archived prior-use
category are unchanged.

Fifth, a documentary review applies only to the filing set it covered. Each
row of \path{site_lot_decisions.csv} lists its \texttt{expected\_component\_jobs},
and the producer checks them against current membership. Charles Place's
agreement assigns complete earlier lots 7 and 71 to the development, but it
does not adjudicate the older filing \texttt{321590541}. Broadway's
9,262-square-foot building parcel likewise leaves its older hotel/residential
filing open. Shepherd/Highland keeps its full 20,653-square-foot earlier
parcel; a 21-unit filing on lot 148 in September 2026 falls after the July 8
cutoff and adds no units to the observed 198-unit pair. At Jamaica/165th
Street, recorded instrument 2026051400252004 separates six residential parcels
totaling 39,349.50 square feet from the 72,982-square-foot MTA terminal. It
does not print DOB job numbers, so its match to the six saved filings is
collective. Housing Database Class~A counts remain the unit measure
throughout, and an economic parent is not thereby a verified legal
wage-assessment unit.

\paragraph{Accepted parent mergers.}
The review accepted seven developments as single economic parents. All seven
are now applied through \path{pair_decisions.csv} and
\path{site_lot_decisions.csv}. Each constituent keeps its administrative units,
original filing date and refiling fields.
\begin{itemize}
\item \textit{Bedford Square} (historical): 2201 Beverly, 2366 Bedford, 2363
Bedford and 158 Lott, filed within 31 days, $296+354+132+95=877$ units. DEC's
remedial investigation describes one four-building redevelopment of the
former Sears parcels, and the developer identifies the four addresses as
buildings A--D. The merger reduces the parent count by three. Ground is the
two surveyed west housing lots (43,413.10 and 73,457.69 square feet) plus the
whole earlier east lot counted once (60,650). Later east areas sum to 60,726;
the 76-square-foot vintage difference is kept as a comparison.
\item \textit{Starhill} (historical): 244- and 326-unit filings five days apart
in February 2020, 570 units. The developer and architect describe one campus.
The March 2022 declaration defines and partitions the combined ground and
excludes retained lots 129 and 162. Later financing phases do not split the
filing-cohort decision.
\item \textit{Godwin/Kimberly} (post-policy): $99+65=164$ units, both filed
March 31, 2025 with the same owner and architect. On September 16 the link was
not adopted: the automatic rule rejects their shared parcel history because
its recorded change postdates filing, and the documents then read did not
allocate the ground between the buildings. The July 2026 agreement then assigned the two buildings to lots 79 and
88 and the existing commercial building to retained lot 78, and the link was
accepted on September 20. Ground is $12,541+5,900=18,441$ square feet. Old lot
78's 17,670-square-foot attribute conflicts with its roughly
28,483-square-foot mapped extent, so it is not allocated proportionally.
\item \textit{GO Broome} (historical): the Suffolk and Norfolk buildings, 35 days
apart, $378+117=495$ units. City Planning identifies one development site with a
shared courtyard on old lots 37 and 75; neighboring lots 1 and 95 are retained.
\item \textit{Woodside} (historical): two filings 13 days apart, 478 units. DEC
names both addresses, one owner and the planned merger of six earlier lots,
whose archival areas sum exactly to the official two-building site.
\item \textit{Motto} (historical): two filings, 264 units. A recorded easement
names both jobs on common developer parcel 38/39/155, matching a June 2019
site plan made before filing. Retained lot 37 is excluded, along with its
building floor.
\item \textit{Elara} (post-policy): $330+99=429$ units filed one day apart in
October 2025. The recorded joint loan names both borrowers, all three parcels
and both exact jobs, and records joint acquisition financing before filing.
\end{itemize}

\paragraph{Relationships rejected or left open.}
Three possible changes to membership were not made. Rockaway Village's three same-day
filings remain one parent; separate financing phases alone do not justify a
split. Fifteen earlier two-family filings at Beach 29/30 have no first permits
and match fifteen retained narrow parcels with their own 30-unit allocation.
They do not imply shared apartment ground; only their documented land
conflicts are cleared, and the filings are preserved. At East 178th Street
only replacement job \texttt{X01273665-I1} is additive; the earlier proposal's
lot 87 lies outside the recorded developer parcel.

The membership review table carries three unresolved pairs.
\textit{Franklin} (\texttt{B01325937-I1}, \texttt{B01325938-I1}) has 259- and
117-unit filings by one developer 22 days apart. A recorded declaration
redraws both lots 63 and 66, and neither parent boundary can be accepted until
the two buildings are shown to be one development or separate decisions.
\textit{Sutphin/95th Avenue} (\texttt{Q00657653}, \texttt{Q00657656}) are
same-day filings with one DOB owner and architect on a prefiling partition of
one 60,138-square-foot lot. Their complete building lots of 30,071 and 30,067
square feet are settled, but the applicant describes neighboring projects, so
both parent IDs are kept pending direct joint-development evidence. At
\textit{1580 Story Avenue}, archived filings \texttt{220700873} and
\texttt{X00756455} each propose 562 units on the same site. The records do not
show whether they are alternatives or phases; neither is merged or deleted.

Other open questions concern ground or filing roles rather than links. Charles
Place's and Broadway's older filings are unresolved. Rockaway III/IV's deed
conveys condominium interests in land shared by several phases, and the
ownership percentages do not measure each phase's physical ground. At 310/322
Grand Concourse the ground evidence is stronger, but the allocation of earlier
zoning or floor is not settled. Atlantic's 52,000-square-foot candidate lacked
a job-specific boundary when last reviewed on September 20.

\paragraph{Adopted development ground.}
The table is built from
\path{tasks/parent_opportunities_manual/output/site_lot_decisions.csv}, which
is the authority for these values. Ground is the sum of
\texttt{development\_area\_sqft} over each parent's rows. Numbers quoted in the
prose above illustrate the rules; the CSV governs any difference.

\begin{center}
\footnotesize
\begin{tabular}{llrll}
\hline
Parent & Sample & Ground (sq ft) & Earlier floor & Ground basis \\
\hline
16 Dupont Street & Hist. & 20,901 & none recorded & Recorded legal area \\
1580 Story Avenue & Hist. & 203,910 & none recorded & Successor area; recorded boundary \\
261/315 Grand Concourse & Hist. & 37,661.3 & kept in full & Recorded survey \\
265 South Street & Hist. & 41,479 & retained excluded & Agency lot reconfiguration \\
270 East 203rd Street & Hist. & 20,000 & kept in full & Archival area \\
355 Exterior Street & Hist. & 120,272 & retained excluded & Successor area; exact-job site plan \\
484 East 178th Street & Hist. & 40,104 & none recorded & Archival area \\
524 Fort Washington Avenue & Hist. & 45,971 & kept in full & Archival area \\
Amsterdam (1440) & Hist. & 24,990 & retained excluded & Successor area; agency land transfer \\
Bedford Square & Hist. & 177,520.79 & retained excluded & Recorded legal and archival areas \\
Casa Celina (Thieriot) & Hist. & 12,613 & retained excluded & Successor area; ground lease \\
Crown Street & Hist. & 55,385 & none recorded & Successor area; exact-job declaration \\
De Kalb Avenue & Hist. & 28,650 & retained excluded & Successor area; development agreement \\
Dean Street (B12/B13) & Hist. & 116,535 & none recorded & Archival area \\
Eagle/West (Site D) & Hist. & 106,018 & none recorded & Recorded survey, reconciled \\
East 125th Street & Hist. & 42,540 & approved estimate & Agency site; survey \\
Flatbush tower & Hist. & 12,603 & approved estimate & Agency ground lease \\
GO Broome & Hist. & 32,395 & archive corrected to zero & Archival area \\
Motto & Hist. & 24,880 & retained excluded & Successor area; prefiling site plan \\
Orchard Street & Hist. & 82,400 & kept in full & Archival area \\
Penn South & Hist. & 40,261 & retained excluded & Successor area; ground lease \\
Shell Road & Hist. & 53,978 & none recorded & Successor area; legal dimensions \\
St.~James/Jerome & Hist. & 17,775 & approved proxy & Successor area; deed and lease \\
Starhill & Hist. & 70,063 & retained excluded & Later area in recorded boundary \\
Van Cortlandt Park South & Hist. & 56,048 & kept in full & Recorded deed area \\
Wallabout/Union & Hist. & 39,323 & none recorded & Successor areas; mapped allocation \\
West 48th Street & Hist. & 21,924 & none recorded & Successor area; agency site \\
Willets Phase I & Hist. & 135,609 & none recorded & Recorded survey \\
Woodside & Hist. & 71,862 & kept in full & Archival area equal to agency site \\
27-30 21st Street & Post & 5,000 & kept in full & Archival area \\
910 Onderdonk Avenue & Post & 9,310 & approved estimate & Successor area; recorded area chart \\
Beach 29/30 & Post & 47,365.13 & none recorded & Recorded development schedule \\
Charles Place & Post & 12,590 & kept in full & Archival area \\
East 108th Street & Post & 44,200 & retained excluded & Recorded legal dimensions \\
East 178th Street & Post & 7,287 & kept in full & Archival area \\
Elara & Post & 39,003 & kept in full & Archival area \\
Godwin/Kimberly & Post & 18,441 & retained excluded & Successor area; legal dimensions \\
Jamaica/165th Street & Post & 39,349.5 & approved estimate & Recorded legal area \\
Kingsbrook/Schenectady & Post & 105,382 & approved estimate & Recorded legal area \\
Noble/Oak & Post & 173,834 & kept in full & City-stated site area \\
North 8th Street & Post & 7,237.85 & retained excluded & Computed from legal courses \\
Ocean Parkway & Post & 20,000 & retained excluded & Recorded legal dimensions \\
Walton/Burnside & Post & 23,000 & kept in full & Archival area \\
\hline
\end{tabular}
\end{center}
\noindent \textit{Earlier floor.} ``None recorded'' means the reference
parcels carry zero building floor. ``Kept in full'' keeps all recorded floor on
the reference parcels. ``Retained excluded'' removes the floor on retained
neighboring buildings (\texttt{excluded\_building\_area\_sqft}). ``Approved
estimate'' and ``approved proxy'' mark the six parents with
\texttt{built\_floor\_area\_estimated = TRUE}; the CSV records that Jacob
approved each flagged allocation or later administrative proxy, and the
excluded amount there is an accounting residual rather than a measured
retained floor. \textit{Ground basis} paraphrases \texttt{area\_basis}.
``Archival area'' is the recorded lot area in the named pre-filing release;
``successor area'' is the later administrative area of the documented
successor parcel.

Several reviewed boundaries needed no production row because the automatic
area was already correct. Sullivan/Empire's March 2026 declaration (ACRIS
2026030900544002, p.~3) partitions old Brooklyn block 1306 lot 28 into the four
tentative lots of its four 99-unit filings, confirming the 38,118-square-foot
production area; the rounded diagram area of 38,119 agrees. Shepherd/Highland,
Broadway and Sutphin/95th are likewise recorded only in the audit's
document-review table.

\paragraph{Remaining flags and the end of per-parent review.}
A parent is flagged when the automatic parcel and filing checks cannot confirm
its land or filing coverage. A flag is an open question, not a confirmed
error, and never removes a parent from the production panels. On September 16
the screen flagged 160 of the 877 weighting-eligible parents. It now flags
145 of 905 (116 historical and 29 post-policy). The two
counts are not a like-for-like reduction: the mergers changed the parent count,
and the September 22 historical-source rule changed the historical sample. On
September 22 the per-parent research batches were not continued after a
land-measurement sensitivity test
(\path{tasks/audits/audit_land_measurement_sensitivity}) showed that the
remaining flagged parents do not move the headline comparison. That test is
described in its own entry.

\textit{Reproduction note.} The adopted decisions are
\path{tasks/parent_opportunities_manual/output/site_lot_decisions.csv} and
\path{pair_decisions.csv}, read by \texttt{build\_parent\_site\_characteristics}
and \texttt{construct\_parent\_cohorts}; root \texttt{make data} rebuilds the
panels. Accepted and open links are in
\path{tasks/audits/audit_parent_site_boundaries/code/parent_site_membership_review.csv},
and reviewed boundaries not needing a production row are in
\path{parent_site_document_review.csv} beside it. The flag counts come from
that audit's \path{output/parent_site_scope_summary.csv}. The documentary
evidence is kept under \path{tasks/audits/parent_review_documents}: each dated
batch folder in \path{output/} holds the committed captures, notes and a
\texttt{sources.sha256} checksum list. The PDFs themselves are outside Git in
\path{data_raw/parent_review_documents/}, and the checksum lists record their
bytes.
